\documentclass[11pt,a4paper]{report}

% ─── Packages ───
\usepackage[utf8]{inputenc}
\usepackage[T1]{fontenc}
\usepackage[margin=1in, top=1.2in, bottom=1.2in]{geometry}
\usepackage{amsmath,amssymb,amsthm,mathtools}
\usepackage{xcolor}
\usepackage{titlesec}
\usepackage{titletoc}
\usepackage{fancyhdr}
\usepackage{enumitem}
\usepackage{tcolorbox}
\usepackage{booktabs}
\usepackage{graphicx}
\usepackage{hyperref}
\usepackage{cleveref}
\usepackage{tikz}
\usepackage{setspace}
\usepackage{parskip}
\usepackage{caption}
\usepackage{etoolbox}
\usepackage{array}
\usepackage{multirow}

\usetikzlibrary{arrows.meta, positioning, shapes.geometric, calc, decorations.pathmorphing, fit, backgrounds, chains, patterns}
\tcbuselibrary{skins,breakable,theorems}

% ─── Colors ───
\definecolor{darknavy}{HTML}{0d1b2a}
\definecolor{accentwhite}{HTML}{e0e1dd}
\definecolor{deepslate}{HTML}{1b263b}
\definecolor{softgold}{HTML}{c8a951}
\definecolor{darkbg}{HTML}{080f1a}
\definecolor{sectionblue}{HTML}{415a77}
\definecolor{subsecblue}{HTML}{778da9}
\definecolor{defcolor}{HTML}{1565c0}
\definecolor{defbg}{HTML}{e3f2fd}
\definecolor{principlecolor}{HTML}{2e7d32}
\definecolor{principlebg}{HTML}{e8f5e9}
\definecolor{mechcolor}{HTML}{00695c}
\definecolor{mechbg}{HTML}{e0f2f1}
\definecolor{mathbg}{HTML}{f8f5ee}
\definecolor{notepurple}{HTML}{6a1b9a}
\definecolor{notebg}{HTML}{f3e5f5}
\definecolor{lightgray}{HTML}{f0f0f0}
\definecolor{darkgray}{HTML}{333333}
\definecolor{histcolor}{HTML}{4e342e}
\definecolor{histbg}{HTML}{efebe9}
\definecolor{conceptcolor}{HTML}{283593}
\definecolor{conceptbg}{HTML}{e8eaf6}
\definecolor{guidecolor}{HTML}{bf360c}
\definecolor{guidebg}{HTML}{fbe9e7}
\definecolor{seriescolor}{HTML}{0d47a1}
\definecolor{seriesbg}{HTML}{e3f2fd}
\definecolor{predcolor}{HTML}{e65100}
\definecolor{objcolor}{HTML}{b71c1c}

% ─── Hyperref Setup ───
\hypersetup{
    colorlinks=true,
    linkcolor=deepslate,
    citecolor=deepslate,
    urlcolor=sectionblue,
    bookmarks=true,
    bookmarksnumbered=true,
    pdfauthor={Comprehensive Guide Series},
    pdftitle={Introduction to Consciousness Studies},
    pdfsubject={Consciousness, Philosophy of Mind, Neuroscience, The Hard Problem},
}

% ─── Header/Footer ───
\pagestyle{fancy}
\fancyhf{}
\fancyhead[L]{\small\color{gray} Introduction to Consciousness Studies}
\fancyhead[R]{\small\color{gray} Comprehensive Reading Material}
\fancyfoot[C]{\thepage}
\renewcommand{\headrulewidth}{0.4pt}
\renewcommand{\footrulewidth}{0pt}
\fancypagestyle{plain}{\fancyhf{}\fancyfoot[C]{\thepage}\renewcommand{\headrulewidth}{0pt}}

% ─── Chapter/Section Formatting ───
\titleformat{\chapter}[display]
  {\normalfont\huge\bfseries\color{deepslate}}
  {\chaptertitlename\ \thechapter}{20pt}{\Huge}
\titlespacing*{\chapter}{0pt}{-20pt}{30pt}
\titleformat{\section}{\normalfont\Large\bfseries\color{sectionblue}}{\thesection}{1em}{}
\titleformat{\subsection}{\normalfont\large\bfseries\color{subsecblue}}{\thesubsection}{1em}{}
\titleformat{\subsubsection}{\normalfont\normalsize\bfseries\color{darkgray}}{\thesubsubsection}{1em}{}

% ─── Tcolorbox environments ───
\newtcolorbox{principlebox}[1][]{enhanced, breakable, colback=principlebg, colframe=principlecolor, coltitle=white, fonttitle=\bfseries, title=#1, left=8pt, right=8pt, top=6pt, bottom=6pt, boxrule=0.5pt, arc=3pt, borderline west={3pt}{0pt}{principlecolor}, before skip=10pt, after skip=10pt}
\newtcolorbox{mechbox}[1][]{enhanced, breakable, colback=mechbg, colframe=mechcolor, coltitle=white, fonttitle=\bfseries, title=#1, left=8pt, right=8pt, top=6pt, bottom=6pt, boxrule=0.5pt, arc=3pt, borderline west={3pt}{0pt}{mechcolor}, before skip=10pt, after skip=10pt}
\newtcolorbox{definitionbox}[1][]{enhanced, breakable, colback=defbg, colframe=defcolor, coltitle=white, fonttitle=\bfseries, title=#1, left=8pt, right=8pt, top=6pt, bottom=6pt, boxrule=0.5pt, arc=3pt, borderline west={3pt}{0pt}{defcolor}, before skip=10pt, after skip=10pt}
\newtcolorbox{mathbox}{enhanced, colback=mathbg, colframe=mathbg!70!black, left=8pt, right=8pt, top=6pt, bottom=6pt, boxrule=0.3pt, arc=2pt, before skip=8pt, after skip=8pt}
\newtcolorbox{notebox}[1][Note]{enhanced, breakable, colback=notebg, colframe=notepurple, coltitle=white, fonttitle=\bfseries, title=#1, left=8pt, right=8pt, top=6pt, bottom=6pt, boxrule=0.5pt, arc=3pt, borderline west={3pt}{0pt}{notepurple}, before skip=10pt, after skip=10pt}
\newtcolorbox{histbox}[1][]{enhanced, breakable, colback=histbg, colframe=histcolor, coltitle=white, fonttitle=\bfseries, title=#1, left=8pt, right=8pt, top=6pt, bottom=6pt, boxrule=0.5pt, arc=3pt, borderline west={3pt}{0pt}{histcolor}, before skip=10pt, after skip=10pt}
\newtcolorbox{conceptbox}[1][]{enhanced, breakable, colback=conceptbg, colframe=conceptcolor, coltitle=white, fonttitle=\bfseries, title=#1, left=8pt, right=8pt, top=6pt, bottom=6pt, boxrule=0.5pt, arc=3pt, borderline west={3pt}{0pt}{conceptcolor}, before skip=10pt, after skip=10pt}
\newtcolorbox{guidebox}[1][]{enhanced, breakable, colback=guidebg, colframe=guidecolor, coltitle=white, fonttitle=\bfseries, title=#1, left=8pt, right=8pt, top=6pt, bottom=6pt, boxrule=0.5pt, arc=3pt, borderline west={3pt}{0pt}{guidecolor}, before skip=10pt, after skip=10pt}
\newtcolorbox{seriesbox}[1][]{enhanced, breakable, colback=seriesbg, colframe=seriescolor, coltitle=white, fonttitle=\bfseries, title=#1, left=8pt, right=8pt, top=6pt, bottom=6pt, boxrule=0.5pt, arc=3pt, borderline west={3pt}{0pt}{seriescolor}, before skip=10pt, after skip=10pt}

\setlength{\parindent}{0pt}
\setlength{\parskip}{6pt}

% ══════════════════════════════════════════════════════════════
\begin{document}

% ─── COVER PAGE ───
\begin{titlepage}
\begin{tikzpicture}[remember picture, overlay]
  \fill[darkbg] (current page.south west) rectangle (current page.north east);
  \fill[softgold] ([yshift=-11cm]current page.north west) rectangle ([yshift=-10.85cm]current page.north east);
  \foreach \a/\r/\op in {0/2.0/0.04, 30/2.3/0.03, 60/1.8/0.05, 90/2.5/0.03, 120/2.1/0.04, 150/1.9/0.04, 180/2.4/0.03, 210/2.0/0.04, 240/2.2/0.03, 270/1.7/0.05, 300/2.3/0.04, 330/2.0/0.03} {
    \draw[accentwhite, opacity=\op, line width=0.4pt] (6, -14.5) -- ({6+\r*cos(\a)}, {-14.5+\r*sin(\a)});
    \fill[accentwhite, opacity=0.06] ({6+\r*cos(\a)}, {-14.5+\r*sin(\a)}) circle (0.08);
  }
  \draw[accentwhite, opacity=0.06, line width=0.5pt] (6, -14.5) circle (1.2);
  \draw[accentwhite, opacity=0.04, line width=0.5pt] (6, -14.5) circle (2.0);
  \draw[accentwhite, opacity=0.03, line width=0.5pt] (6, -14.5) circle (2.8);
  \fill[softgold, opacity=0.08] (6, -14.5) circle (0.2);
\end{tikzpicture}

\vspace*{4cm}
\begin{center}
  {\fontsize{36}{44}\selectfont\bfseries\color{white} Introduction to\\[8pt] Consciousness Studies}\\[24pt]
  {\Large\color{gray!70} Foundations, Problems, and the Landscape of Theories}\\[20pt]
  {\color{softgold}\rule{6cm}{1pt}}\\[20pt]
  {\large\color{gray!60} A Comprehensive Guide to the Science and Philosophy\\[4pt] of Subjective Experience}\\[40pt]
  {\color{softgold}\large Volume 0 of the Consciousness Theory Series}\\[8pt]
  {\color{gray!50}\normalsize The Hard Problem $\bullet$ Neural Correlates $\bullet$ Qualia $\bullet$ Methodology}\\[60pt]
  {\color{gray!40}\small Comprehensive Study Material \quad$\vert$\quad February 2026}
\end{center}
\end{titlepage}

% ─── TABLE OF CONTENTS ───
\tableofcontents
\thispagestyle{plain}
\newpage


% ══════════════════════════════════════════════════════════════
\chapter{What is Consciousness?}
% ══════════════════════════════════════════════════════════════

\section{The Phenomenon}

You are reading this sentence. You see marks on a page or screen. You understand the words. You are aware of understanding them. There is \emph{something it is like} to be you right now---a felt quality to your experience, a subjective perspective on the world. This is \textbf{consciousness}.

Consciousness is at once the most familiar thing in the universe---it is the only thing any of us knows directly---and the most puzzling. We know what it is from the inside but struggle to explain it from the outside. As Thomas Nagel (1974) framed it: a being is conscious if and only if \textbf{there is something it is like to be that being}.

\begin{definitionbox}[Working Definitions]
\textbf{Consciousness} (phenomenal consciousness): The subjective, qualitative character of experience. The ``what it is like-ness'' of mental states. When you see red, feel pain, or taste coffee, there is a phenomenal quality---a \emph{quale} (plural: \emph{qualia})---that constitutes your experience.

\textbf{Access consciousness} (Block, 1995): A mental state is access-conscious if its content is available for verbal report, reasoning, and the rational control of behavior. Access consciousness concerns the \emph{functional} role of a mental state, not its subjective quality.

\textbf{Self-consciousness:} Awareness of oneself as a subject of experience. A higher-order capacity that involves modeling one's own mental states.

The central mystery concerns \textbf{phenomenal consciousness}: why and how do physical processes in the brain give rise to subjective experience?
\end{definitionbox}

\section{Why Consciousness is Hard to Study}

\begin{conceptbox}[The Unique Challenges of Consciousness]
\begin{enumerate}[leftmargin=1.5em]
  \item \textbf{First-person vs.\ third-person:} All other scientific phenomena are observable from the outside (third-person). Consciousness is known directly only from the inside (first-person). There is no ``consciousness-meter''; we can only measure neural activity and infer consciousness.
  \item \textbf{The explanatory gap:} Even a complete physical description of the brain seems to leave out something---the \emph{felt quality} of experience.
  \item \textbf{Privacy:} Your consciousness is directly accessible only to you. I can observe your behavior and brain activity but can never directly access your experience.
  \item \textbf{Ambiguity of terms:} ``Consciousness'' is used to mean many different things---wakefulness, awareness, attention, self-reflection, sentience---creating confusion across disciplines.
\end{enumerate}
\end{conceptbox}


% ══════════════════════════════════════════════════════════════
\chapter{The Problems of Consciousness}
% ══════════════════════════════════════════════════════════════

\section{The Hard Problem}

\begin{definitionbox}[The Hard Problem of Consciousness (Chalmers, 1995)]
\textbf{The Hard Problem:} Why is there subjective experience at all? Why don't all the physical processes in the brain go on ``in the dark,'' without any accompanying inner experience? Even if we explain every functional aspect of consciousness, there would remain a further question: \emph{why is there something it is like} to undergo these processes?

Contrast with the \textbf{Easy Problems} (``easy'' only relative to the Hard Problem):
\begin{itemize}[leftmargin=1.5em]
  \item How does the brain integrate information from different senses?
  \item How does attention select certain inputs for processing?
  \item How can we report on our mental states?
  \item What distinguishes wakefulness from sleep?
\end{itemize}
These ask how the brain \emph{does} things. The Hard Problem asks why doing those things \emph{feels like} anything.
\end{definitionbox}

\subsection{The Zombie Argument}

\begin{conceptbox}[The Zombie Argument]
A \textbf{philosophical zombie} is a being physically identical to a conscious human---same neurons, same connections, same firing patterns---but with \emph{no subjective experience}. It behaves identically but ``the lights are off inside.''

\begin{enumerate}[leftmargin=1.5em]
  \item Zombies are \emph{conceivable}: we can coherently imagine a physical duplicate without consciousness.
  \item If zombies are conceivable, consciousness is not logically entailed by physical properties.
  \item Therefore, consciousness is something \emph{over and above} the physical.
\end{enumerate}
This argument is controversial---many physicalists deny that zombies are genuinely conceivable---but it sharpens the challenge.
\end{conceptbox}

\section{The Explanatory Gap}

The \textbf{explanatory gap} (Levine, 1983) is the observation that even a complete neuroscientific account of the brain leaves a gap between the physical story and phenomenal reality. We can explain \emph{why} water is H$_2$O. But we cannot explain \emph{why} certain neural activity is accompanied by the experience of red rather than green---or by any experience at all.

\section{The Binding Problem}

\begin{definitionbox}[The Binding Problem]
Conscious experience is \textbf{unified}: you see a single visual scene with shape, color, motion, and depth all bound together. But the brain processes these features in anatomically separate regions. How are they combined?

Sub-problems: feature binding (red + round + moving for a single object), object binding (separating objects in a scene), cross-modal binding (sight + sound + touch), and temporal binding (continuity over time).
\end{definitionbox}

\section{The Neural Correlates of Consciousness (NCC)}

\begin{definitionbox}[Neural Correlates of Consciousness]
The \textbf{NCC} is the minimal set of neuronal events and mechanisms jointly sufficient for a specific conscious experience (Koch et al., 2016).
\begin{itemize}[leftmargin=1.5em]
  \item \textbf{Content NCC:} The neural basis of a \emph{specific} experience (seeing red, hearing a tone).
  \item \textbf{Full NCC:} The neural basis of consciousness itself (what distinguishes conscious from unconscious states).
\end{itemize}
Finding NCCs is the central empirical project of consciousness science. But NCCs alone do not solve the Hard Problem: they tell us \emph{which} neural events correlate with consciousness, not \emph{why}.
\end{definitionbox}


% ══════════════════════════════════════════════════════════════
\chapter{Key Concepts and Distinctions}
% ══════════════════════════════════════════════════════════════

\section{Levels and Contents of Consciousness}

\begin{principlebox}[Two Dimensions of Consciousness]
\textbf{Level} (arousal, wakefulness): A graded dimension from deep coma through sleep, sedation, drowsiness, to full alert wakefulness. Measured clinically by the Glasgow Coma Scale, and experimentally by the perturbational complexity index (PCI).

\textbf{Content} (what you are experiencing): The specific qualitative character of experience at a given moment. Content varies independently of level: you can be fully awake with sparse content (sensory deprivation) or have rich content at reduced levels (dreaming).
\end{principlebox}

\section{Access vs.\ Phenomenal Consciousness}

This distinction (Block, 1995) is central to debates between theories:

\begin{center}
\renewcommand{\arraystretch}{1.4}
\begin{tabular}{>{\raggedright}p{3.5cm} >{\raggedright}p{5cm} >{\raggedright\arraybackslash}p{5cm}}
  \toprule
  & \textbf{Phenomenal Consciousness} & \textbf{Access Consciousness} \\
  \midrule
  Definition & Subjective, qualitative character & Availability for report, reasoning, action \\
  Key question & What is it \emph{like}? & What can you \emph{do} with it? \\
  Measurable? & Indirectly (through report) & Directly (through behavior) \\
  Can dissociate? & Possibly (phenomenal overflow?) & Possibly (blindsight?) \\
  Key theories & RPT, IIT (phenomenal is primary) & GWT, HOT (access \emph{is} consciousness) \\
  \bottomrule
\end{tabular}
\end{center}

Whether phenomenal consciousness can exist without access consciousness is one of the most contested questions. If it can (``phenomenal overflow''), this favors RPT and IIT. If not, this favors GWT and HOT.

\section{Qualia}

\begin{definitionbox}[Qualia]
\textbf{Qualia} (singular: \emph{quale}) are the subjective, qualitative properties of experience: the redness of red, the painfulness of pain. Key features:
\begin{itemize}[leftmargin=1.5em]
  \item \textbf{Intrinsic:} A quale is a property of the experience itself.
  \item \textbf{Private:} Only the subject can know their own qualia directly.
  \item \textbf{Ineffable:} Qualia cannot be fully communicated to someone who hasn't experienced them.
  \item \textbf{Directly apprehended:} You know your qualia simply by having them.
\end{itemize}
Whether qualia exist as described is itself debated. Dennett (1988) argued against qualia; others (Chalmers, Block) consider them central.
\end{definitionbox}

\section{Key Thought Experiments}

\begin{histbox}[Foundational Thought Experiments]
\textbf{Nagel's Bat (1974):} We cannot know what it is like to be a bat, no matter how much we know about bat neuroscience. This illustrates the irreducibility of subjective experience.

\textbf{Jackson's Mary (1982):} Mary knows everything physical about color vision but has lived in a black-and-white room. When she sees red for the first time, does she learn something new? If yes, physical knowledge is incomplete.

\textbf{Chalmers' Zombies (1996):} A being physically identical to you but without consciousness. If conceivable, consciousness is not logically entailed by physics.

\textbf{Searle's Chinese Room (1980):} A person follows rules to manipulate Chinese symbols without understanding Chinese. Challenges computational theories: syntax doesn't produce semantics.

\textbf{Block's China Brain (1978):} If the entire population of China simulated a brain, would the system be conscious? Tests intuitions about substrate-independence.
\end{histbox}


% ══════════════════════════════════════════════════════════════
\chapter{The Philosophical Landscape}
% ══════════════════════════════════════════════════════════════

\section{Major Positions}

\begin{center}
\renewcommand{\arraystretch}{1.3}
\begin{tabular}{>{\raggedright}p{3.5cm} >{\raggedright}p{4cm} >{\raggedright\arraybackslash}p{6cm}}
  \toprule
  \textbf{Position} & \textbf{Core Claim} & \textbf{Key Proponents} \\
  \midrule
  Physicalism & Consciousness is entirely physical & Dennett, Churchland, Crick, Koch \\
  Property dualism & Non-physical property of physical systems & Chalmers \\
  Substance dualism & Mind and matter are distinct substances & Descartes \\
  Functionalism & Consciousness defined by functional role & Putnam, Fodor \\
  Illusionism & Phenomenal consciousness is an illusion & Dennett, Frankish \\
  Panpsychism & Consciousness is fundamental to all matter & Goff, Strawson \\
  Idealism & Consciousness is the only fundamental reality & Kastrup, Hoffman \\
  Neutral monism & Mind and matter are aspects of neutral stuff & James, Russell \\
  Mysterianism & Humans cannot solve the problem & McGinn \\
  \bottomrule
\end{tabular}
\end{center}

\section{Physicalism and Its Varieties}

\begin{itemize}[leftmargin=1.5em]
  \item \textbf{Reductive physicalism (identity theory):} Conscious states \emph{are} brain states. Pain just \emph{is} C-fiber firing.
  \item \textbf{Functionalism:} Conscious states are defined by functional roles, not substrate. Allows multiple realizability: consciousness in silicon, aliens, etc.
  \item \textbf{Eliminativism:} Folk-psychological concepts of consciousness are mistaken and will be eliminated by mature neuroscience.
  \item \textbf{Illusionism:} Phenomenal consciousness doesn't exist as such; we merely have functional states that represent themselves as having qualitative character.
\end{itemize}

\section{Dualism}

\begin{itemize}[leftmargin=1.5em]
  \item \textbf{Substance dualism} (Descartes): Mind and body are different substances. Faces the interaction problem.
  \item \textbf{Property dualism} (Chalmers): One substance (physical), two kinds of properties (physical and phenomenal). Phenomenal properties supervene on but are not reducible to physical ones.
  \item \textbf{Epiphenomenalism:} Consciousness exists but has no causal power. Physical events cause conscious events, not vice versa.
\end{itemize}


% ══════════════════════════════════════════════════════════════
\chapter{Experimental Methods}
% ══════════════════════════════════════════════════════════════

\section{Contrastive Methods}

The primary experimental strategy is \textbf{contrastive analysis}: compare brain activity when a stimulus is consciously perceived vs.\ when it is not.

\begin{mechbox}[Key Experimental Paradigms]
\begin{enumerate}[leftmargin=1.5em]
  \item \textbf{Binocular rivalry:} Different images to each eye; perception alternates.
  \item \textbf{Visual masking:} A target is rendered invisible by a subsequent mask.
  \item \textbf{Inattentional blindness:} Subjects fail to notice unexpected stimuli.
  \item \textbf{Change blindness:} Subjects fail to detect large changes across interruptions.
  \item \textbf{Attentional blink:} The second of two rapid targets is often missed.
  \item \textbf{Global-local paradigms:} Stimuli near the threshold of consciousness.
  \item \textbf{Sleep/anesthesia/coma:} Comparing activity across consciousness levels.
  \item \textbf{Perturbational approaches:} TMS to probe connectivity and complexity (PCI).
\end{enumerate}
\end{mechbox}

\section{Neural Measures}

\begin{center}
\renewcommand{\arraystretch}{1.3}
\begin{tabular}{>{\raggedright}p{2.5cm} >{\raggedright}p{3cm} >{\raggedright}p{3cm} >{\raggedright\arraybackslash}p{4cm}}
  \toprule
  \textbf{Method} & \textbf{Spatial Res.} & \textbf{Temporal Res.} & \textbf{What it Measures} \\
  \midrule
  EEG & $\sim$1 cm & $\sim$1 ms & Scalp electric field \\
  MEG & $\sim$5 mm & $\sim$1 ms & Magnetic field \\
  fMRI & $\sim$1 mm & $\sim$1 s & Blood oxygenation (BOLD) \\
  ECoG & $\sim$1 mm & $\sim$1 ms & Cortical surface field \\
  Single-unit & $\sim$10 $\mu$m & $\sim$0.1 ms & Individual neuron firing \\
  TMS & Variable & $\sim$1 ms & Causal perturbation \\
  \bottomrule
\end{tabular}
\end{center}

\section{The Adversarial Collaboration Framework}

\begin{principlebox}[Adversarial Collaborations in Consciousness Science]
The Templeton World Charity Foundation funded large-scale adversarial collaborations testing predictions of major theories (initially IIT vs.\ GWT, later expanded). Key features:
\begin{itemize}[leftmargin=1.5em]
  \item Pre-registered predictions from each theory's proponents
  \item Agreed-upon experimental protocols
  \item Large sample sizes and multi-site replication
  \item Independent data analysis
  \item Results published regardless of outcome
\end{itemize}
Initial results (2023) found mixed evidence: neither IIT nor GWT was fully supported. The sustained posterior negativity predicted by IIT was partially confirmed; the prefrontal ignition predicted by GWT received weaker support. These results have reshaped the field, showing that the truth may involve elements of multiple theories.
\end{principlebox}


% ══════════════════════════════════════════════════════════════
\chapter{A Brief History of Consciousness Science}
% ══════════════════════════════════════════════════════════════

\section{Ancient and Early Modern}

Consciousness has been a central concern of philosophy since antiquity. Plato distinguished the soul from the body; Aristotle explored perception and thought; Descartes (1641) made consciousness the foundation of certainty (``\emph{Cogito, ergo sum}'') and established the mind-body problem in its modern form.

\section{The 19th Century}

William James' \textit{The Principles of Psychology} (1890) provided the first systematic treatment of consciousness from a scientific perspective, introducing the ``stream of consciousness'' and exploring attention, perception, and selfhood. Wilhelm Wundt founded the first experimental psychology laboratory (1879).

\section{The Behaviorist Exile (1920s--1960s)}

Behaviorism (Watson, Skinner) banished consciousness from scientific psychology. Only observable behavior was considered scientific. This lasted roughly four decades.

\section{The Cognitive Revolution (1950s--1970s)}

The cognitive revolution (Chomsky, Miller, Newell, Simon) restored internal mental states to scientific respectability, but focused on \emph{information processing}---not subjective experience.

\section{The Modern Era (1990s--Present)}

The scientific study of consciousness became respectable in the 1990s, driven by:
\begin{itemize}[leftmargin=1.5em]
  \item \textbf{Crick \& Koch (1990):} Proposed that consciousness could be studied by finding its neural correlates.
  \item \textbf{Chalmers (1995):} Formulated the Hard Problem.
  \item \textbf{Baars (1988):} Proposed Global Workspace Theory.
  \item \textbf{Tononi (2004):} Proposed Integrated Information Theory.
  \item \textbf{Advances in neuroimaging:} fMRI, EEG/MEG advances, optogenetics.
  \item \textbf{Adversarial Collaborations (2019--):} Systematic tests of competing theories.
\end{itemize}


% ══════════════════════════════════════════════════════════════
\chapter{The Landscape of Theories}
% ══════════════════════════════════════════════════════════════

\section{Taxonomy}

\begin{principlebox}[Organizing Dimensions]
\textbf{By level of explanation:}
\begin{itemize}[leftmargin=1.5em]
  \item \textbf{Neuroscientific:} GWT, RPT, CEMI
  \item \textbf{Mathematical/formal:} IIT, PP
  \item \textbf{Cognitive/representational:} HOT, AST
  \item \textbf{Physical:} Orch OR, CEMI
  \item \textbf{Metaphysical:} Panpsychism
\end{itemize}

\textbf{By stance on the Hard Problem:}
\begin{itemize}[leftmargin=1.5em]
  \item \textbf{Dissolve it:} AST, HOT
  \item \textbf{Address it:} IIT, Panpsychism, Orch OR
  \item \textbf{Set it aside:} GWT, RPT, PP, CEMI
\end{itemize}

\textbf{By substrate commitment:}
\begin{itemize}[leftmargin=1.5em]
  \item \textbf{Substrate-independent:} GWT, HOT, AST, PP
  \item \textbf{Substrate-dependent:} Orch OR, CEMI, IIT
  \item \textbf{Universal:} Panpsychism
\end{itemize}
\end{principlebox}

\section{The Series}

\begin{seriesbox}[The Consciousness Theory Series]
This introduction is the first of 11 volumes. Each provides a comprehensive, self-contained guide to a specific theory. The series:

\begin{center}
\renewcommand{\arraystretch}{1.35}
\begin{tabular}{>{\raggedright}p{0.5cm} >{\raggedright}p{5.5cm} >{\raggedright\arraybackslash}p{7cm}}
  \toprule
  & \textbf{Document} & \textbf{Key Focus} \\
  \midrule
  0 & \textbf{Introduction} (this document) & Foundations, problems, methods, landscape \\
  1 & \textbf{IIT} --- Integrated Information Theory & $\Phi$, axioms, qualia space, mathematical formalism \\
  2 & \textbf{GWT} --- Global Workspace Theory & Ignition, broadcast, thalamocortical dynamics \\
  3 & \textbf{RPT} --- Recurrent Processing Theory & Recurrence, laminar circuits, phenomenal overflow \\
  4 & \textbf{HOT} --- Higher-Order Theories & Meta-representation, signal detection, PFC \\
  5 & \textbf{AST} --- Attention Schema Theory & Attention modeling, schema, TPJ, social cognition \\
  6 & \textbf{PP} --- Predictive Processing & Free energy, prediction error, generative models \\
  7 & \textbf{Orch OR} --- Objective Reduction & Quantum gravity, microtubules, non-computability \\
  8 & \textbf{Panpsychism \& Cosmopsychism} & Fundamental mentality, combination problem \\
  9 & \textbf{CEMI} --- EM Field Theories & Maxwell's equations, ephaptic coupling \\
  10 & \textbf{Master Comparison \& Synthesis} & Cross-theory tables, open frontiers \\
  \bottomrule
\end{tabular}
\end{center}
\end{seriesbox}

\section{Brief Previews}

\subsection{Integrated Information Theory (IIT)}
IIT (Tononi) starts from phenomenological axioms and derives a mathematical framework. Consciousness \emph{is} integrated information ($\Phi$). Has panpsychist implications.

\subsection{Global Workspace Theory (GWT)}
GWT (Baars; Dehaene \& Changeux) proposes that consciousness arises when information is ``broadcast'' to a global workspace, making it available to multiple cognitive processes simultaneously.

\subsection{Recurrent Processing Theory (RPT)}
RPT (Lamme) identifies consciousness with \textbf{recurrent} (feedback) processing in sensory cortex. Predicts phenomenal overflow and locates the NCC in sensory cortex, not prefrontal regions.

\subsection{Higher-Order Theories (HOT)}
HOT (Rosenthal; Brown, Lau, LeDoux) proposes that a mental state is conscious when it is the object of a \textbf{higher-order representation}. Locates NCC in prefrontal cortex.

\subsection{Attention Schema Theory (AST)}
AST (Graziano) proposes consciousness is the brain's internal model of its own attention---an ``attention schema'' analogous to the body schema.

\subsection{Predictive Processing (PP)}
PP (Clark; Friston; Seth) models the brain as a hierarchical prediction machine. Consciousness is a ``controlled hallucination.'' The Free Energy Principle provides a unifying framework.

\subsection{Orchestrated Objective Reduction (Orch OR)}
Orch OR (Penrose \& Hameroff) proposes quantum computations in microtubules terminated by gravitational state reduction. Consciousness is non-computable.

\subsection{Panpsychism \& Cosmopsychism}
Panpsychism (Goff, Strawson, Chalmers) holds that consciousness is fundamental to all matter. A metaphysical framework addressing the Hard Problem. Faces the combination problem.

\subsection{CEMI / EM Field Theories}
CEMI (McFadden) proposes consciousness is the brain's electromagnetic field, which integrates information through superposition and influences neurons through ephaptic coupling.


% ══════════════════════════════════════════════════════════════
\chapter{How to Read This Series}
% ══════════════════════════════════════════════════════════════

\section{Prerequisites}

The guides assume basic neuroscience (neurons, synapses, brain anatomy), undergraduate-level mathematics (calculus, linear algebra, probability, information theory), and basic philosophy of mind (mind-body problem, physicalism, dualism, functionalism---this introduction provides the grounding). Advanced mathematics is introduced within each guide as needed.

\section{Reading Order}

The volumes are self-contained and can be read in any order. Recommended paths:

\begin{guidebox}[Recommended Reading Paths]
\textbf{Neuroscience first:} Introduction $\to$ GWT $\to$ RPT $\to$ IIT $\to$ CEMI $\to$ HOT $\to$ AST $\to$ PP $\to$ Orch OR $\to$ Panpsychism $\to$ Synthesis

\textbf{Philosophy first:} Introduction $\to$ Panpsychism $\to$ IIT $\to$ HOT $\to$ AST $\to$ GWT $\to$ RPT $\to$ PP $\to$ CEMI $\to$ Orch OR $\to$ Synthesis

\textbf{Adversarial collaborations:} Introduction $\to$ IIT $\to$ GWT $\to$ RPT $\to$ HOT $\to$ Synthesis (core four) $\to$ remaining volumes

\textbf{Overview then deep-dive:} Introduction $\to$ Synthesis $\to$ whichever individual volumes interest you most
\end{guidebox}

\section{Conventions}

Each theory guide follows a consistent structure: Introduction $\to$ Prerequisites $\to$ Core Theory $\to$ Neural Implementation $\to$ Experimental Evidence $\to$ Predictions $\to$ Comparisons $\to$ Criticisms $\to$ References.

Visual conventions: \textcolor{defcolor}{blue boxes} for definitions, \textcolor{principlecolor}{green boxes} for key principles, \textcolor{mechcolor}{teal boxes} for mechanisms, \textcolor{notepurple}{purple boxes} for important notes, \textcolor{predcolor}{orange boxes} for predictions/evidence, and \textcolor{objcolor}{red boxes} for objections. Mathematical content appears in tan-shaded boxes.


% ══════════════════════════════════════════════════════════════
\chapter{Further Reading}
% ══════════════════════════════════════════════════════════════

\section*{Essential Introductory Texts}
\begin{enumerate}[label={[\arabic*]}, leftmargin=2.5em]
  \item Seth, A.\ K.\ (2021). \textit{Being You: A New Science of Consciousness}. Faber \& Faber.
  \item Koch, C.\ (2019). \textit{The Feeling of Life Itself}. MIT Press.
  \item Goff, P.\ (2019). \textit{Galileo's Error}. Pantheon.
  \item Dehaene, S.\ (2014). \textit{Consciousness and the Brain}. Viking.
  \item Chalmers, D.\ J.\ (1996). \textit{The Conscious Mind}. Oxford University Press.
\end{enumerate}

\section*{Key Review Articles}
\begin{enumerate}[label={[\arabic*]}, leftmargin=2.5em, start=6]
  \item Seth, A.\ K.\ \& Bayne, T.\ (2022). Theories of consciousness. \textit{Nature Reviews Neuroscience}, 23, 439--452.
  \item Melloni, L., et al.\ (2021). Making the hard problem of consciousness easier. \textit{Science}, 372(6545), 911--912.
  \item Koch, C., et al.\ (2016). Neural correlates of consciousness: progress and problems. \textit{Nature Reviews Neuroscience}, 17, 307--321.
  \item Mashour, G.\ A., et al.\ (2020). Conscious processing and the global neuronal workspace hypothesis. \textit{Neuron}, 105(5), 776--798.
  \item Block, N.\ (2005). Two neural correlates of consciousness. \textit{Trends in Cognitive Sciences}, 9(2), 46--52.
\end{enumerate}

\section*{Foundational Philosophical Works}
\begin{enumerate}[label={[\arabic*]}, leftmargin=2.5em, start=11]
  \item Nagel, T.\ (1974). What is it like to be a bat? \textit{Philosophical Review}, 83(4), 435--450.
  \item Jackson, F.\ (1982). Epiphenomenal qualia. \textit{Philosophical Quarterly}, 32(127), 127--136.
  \item Chalmers, D.\ J.\ (1995). Facing up to the problem of consciousness. \textit{Journal of Consciousness Studies}, 2(3), 200--219.
  \item Levine, J.\ (1983). Materialism and qualia: the explanatory gap. \textit{Pacific Philosophical Quarterly}, 64(4), 354--361.
  \item Block, N.\ (1995). On a confusion about a function of consciousness. \textit{Behavioral and Brain Sciences}, 18(2), 227--247.
  \item Dennett, D.\ C.\ (1991). \textit{Consciousness Explained}. Little, Brown.
  \item Searle, J.\ (1980). Minds, brains, and programs. \textit{Behavioral and Brain Sciences}, 3(3), 417--424.
\end{enumerate}

\section*{Empirical Foundations}
\begin{enumerate}[label={[\arabic*]}, leftmargin=2.5em, start=18]
  \item Crick, F.\ \& Koch, C.\ (1990). Towards a neurobiological theory of consciousness. \textit{Seminars in the Neurosciences}, 2, 263--275.
  \item Dehaene, S.\ \& Changeux, J.-P.\ (2011). Experimental and theoretical approaches to conscious processing. \textit{Neuron}, 70(2), 200--227.
  \item Lamme, V.\ A.\ F.\ (2006). Towards a true neural stance on consciousness. \textit{Trends in Cognitive Sciences}, 10(11), 494--501.
  \item Casali, A.\ G., et al.\ (2013). A theoretically based index of consciousness. \textit{Science Translational Medicine}, 5(198), 198ra105.
  \item Tononi, G.\ \& Koch, C.\ (2015). Consciousness: here, there and everywhere? \textit{Phil.\ Trans.\ Roy.\ Soc.\ B}, 370(1668), 20140167.
\end{enumerate}

\vfill
\begin{center}
  {\color{gray}\rule{0.4\textwidth}{0.5pt}}\\[8pt]
  {\small\color{gray}\itshape End of Introduction --- Continue to the individual theory guides}
\end{center}

\end{document}
